\documentclass[a4paper, 12pt]{article}

\usepackage{utils}

\renewcommand*{\today}{10 February 2026}

\begin{document}

\hotbox{Probabilités 1}{CM 5}{\today}

\begin{definition}
    Soit $X$ une variable aléatoire discrète sur l'espace probabilisé $(\Omega, \mathcal{T}, \mathbb{P})$.
    On dit que $X$ possède \textbf{un moment d'ordre 2} si $\mathbb{E}(X^2)$ existe.

    On appelle \textbf{variance} de $X$ la quantité $\mathrm{Var}(X) = \mathbb{E}[(X - \mathbb{E}[X])^2]$.
\end{definition}

\begin{proposition}{}{}
    Si $\mathbb{E}[X^2]$ existe, alors $\mathbb{E}[X]$ existe également.
\end{proposition}

\begin{demonstration}
    À faire.
\end{demonstration}

\begin{proposition}{}{}
    Soit $X$ une variable aléatoire discrète sur l'espace probabilisé $(\Omega, \mathcal{T}, \mathbb{P})$.
    Si $X$ possède une variance, alors pour tout $a, b \in \mathbb{R}$, $aX + b$ possède une variance et
    $$
    \mathrm{Var}(aX + b) = a^2 \mathrm{Var}(X)
    $$
\end{proposition}

\begin{demonstration}
    $\mathbb{E}[aX + b] = a \mathbb{E}[X] + b$ et
    \begin{align*}
        \mathrm{Var}(aX + b) &= \sum_{i \in I}\left( a x_i + b - (a \mathbb{E}[X] + b) \right)^2 \mathbb{P}(\{ X = x_i \}) \\
        &= a^2 \sum_{i \in I} (x_i - \mathbb{E}[X])^2 \mathbb{P}(\{ X = x_i \}) \\
        &= a^2 \mathrm{Var}(X) 
    \end{align*}
\end{demonstration}

\begin{proposition}{Formule de König-Huygens}{}
    Soit $X$ une variable aléatoire discrète sur l'espace probabilisé $(\Omega, \mathcal{T}, \mathbb{P})$.
    Si $X$ possède une variance, alors
    $$
    \mathrm{Var}(X) = \mathbb{E}[X^2] - (\mathbb{E}[X])^2
    $$
\end{proposition}

\begin{demonstration}
    À faire.
\end{demonstration}

\begin{proposition}{}{}
    Soit $X$ une variable aléatoire discrète sur l'espace probabilisé $(\Omega, \mathcal{T}, \mathbb{P})$.
    Si $X$ possède une variance, alors $\mathrm{Var}(X) = 0 \iff X \equiv \mathbb{E}[X]$.
\end{proposition}

\begin{demonstration}
\end{demonstration}

\end{document}